\section{A metric and volume bound for connected supports}
\label{sec:metric}

We work with the graph \(Y\), group \(H\), and nonempty
\(n\)-dimensional complex \(K\) of Theorem~\ref{thm:sparse}, where
\(n\geq1\).  We need two kinds of geometric control: a metric with
finite neighborhoods in which to make cuts, and Euclidean area
functionals for which those cuts satisfy a slicing inequality.
We will use a scaled transportation metric on \(|K|\) and,
separately inside each top simplex, Euclidean structures determined
by its spanning trees.  Each top simplex will have volume \(1/n!\),
while every affine branch of a distance function will have gradient
norm at most one in each of these structures.  Counting the possible
connected vertex supports will then give the ball-volume estimate
used in the separating argument.

\subsection{The metric and its finite local geometry}

Represent a point of \(|K|\) by a probability vector \(p\) on \(Y\) whose
support lies in a simplex of \(K\).  Let \(d_Y\) denote graph distance and
define
\begin{equation}
 \rho(p,q)=\frac1{\sqrt n}
 \sup_{\operatorname{Lip}_{d_Y}(\ell)\leq1}
 \sum_{v\in Y}\ell(v)\bigl(p(v)-q(v)\bigr).
 \label{eq:transport-metric}
\end{equation}
This is the scaled Kantorovich--Rubinstein metric on finitely
supported probability measures, in its Lipschitz-potential form
\cite[pp.~1381--1382]{Kantorovich2006}. We establish its needed
properties directly.
Only finitely many terms occur in the pairing.  To see that the
supremum is finite and describe the distance function on each simplex,
let \(V\subset Y\) be a nonempty finite set and put
\[
 \mathcal P_V=
 \left\{\ell\in\R^V:
   \sum_{v\in V}\ell(v)=0,\quad
   |\ell(v)-\ell(w)|\leq d_Y(v,w)\quad(v,w\in V)\right\}.
\]
This is a nonempty compact polytope.  Indeed, its difference constraints
and zero mean bound every coordinate by
\(\max_{v,w\in V}d_Y(v,w)\).  Each \(\ell\in\mathcal P_V\) extends to
a \(1\)-Lipschitz function on \(Y\) by
\[
 \widetilde\ell(y)=\min_{v\in V}\bigl(\ell(v)+d_Y(v,y)\bigr).
\]
The difference constraints show that this extension agrees with
\(\ell\) on \(V\), and the triangle inequality proves that it is
\(1\)-Lipschitz.  Conversely, a global \(1\)-Lipschitz function restricts,
after subtraction of its mean, to an element of \(\mathcal P_V\).
Constants pair to zero with \(p-q\).  Consequently, when \(V\) contains
both supports, the supremum in \eqref{eq:transport-metric} is the maximum
over \(\mathcal P_V\).

Symmetry follows by replacing \(\ell\) with \(-\ell\), and the triangle
inequality follows by applying each potential to
\(p-q=(p-u)+(u-q)\).  The indicator of a single vertex is
\(1\)-Lipschitz for \(d_Y\), so
\begin{equation}
 |p(v)-q(v)|\leq\sqrt n\,\rho(p,q)
 \qquad(v\in Y).
 \label{eq:coordinate-control}
\end{equation}
Thus \(\rho\) separates points and is a metric.  It is \(H\)-invariant
because graph automorphisms permute the admissible potentials.  On
probability vectors supported in a fixed finite set, it is the
restriction of a norm on the finite-dimensional space of vectors with
coordinate sum zero.  It therefore gives the usual topology on every
finite subcomplex.

For \(c\in|K|\), write
\[
 \begin{aligned}
 f_c(p)&=\rho(c,p),& U_s(c)&=\{p:f_c(p)<s\},\\
 L_s(c)&=\{p:f_c(p)\leq s\},& S_s(c)&=\{p:f_c(p)=s\}.
 \end{aligned}
\]
On a fixed closed simplex \(\sigma\), the function \(f_c\) is a maximum
of finitely many affine functions: use the extreme points of
\(\mathcal P_{V(\sigma)\cup\supp c}\).  Here \(V(\sigma)\) denotes the
vertex set of \(\sigma\).  Taking all these extreme points makes the
collection of affine branches equivariant under \(H\).

\begin{lemma}\label{lem:proper}
Every closed \(\rho\)-ball meets only finitely many simplices of \(K\).
The complex \(K\) is locally finite, and \(\rho\) gives its usual
realization topology and is proper.  For each \(c\in|K|\), \(s\geq0\),
and simplex \(\kappa\), the set
\[
 \{h\in H:L_s(hc)\cap\kappa\ne\varnothing\}
\]
is finite.  The setwise stabilizer \(H_\kappa\) of a nonempty simplex
is finite as well.
\end{lemma}

\begin{proof}
Let \(A=\supp c\).  The function
\(d_Y(\,\cdot\,,A)\) is \(1\)-Lipschitz and vanishes on \(A\).
Using its negative in \eqref{eq:transport-metric} gives
\begin{equation}
 \min_{\substack{a\in\supp c\\v\in\supp p}}d_Y(a,v)
 \leq \sum_v p(v)d_Y(v,A)
 \leq\sqrt n\,\rho(c,p).
 \label{eq:support-distance}
\end{equation}
If a top simplex \(\sigma\) meets \(L_s(c)\) at \(p\), it therefore
contains a vertex within distance \(\lfloor s\sqrt n\rfloor\) of \(A\).
The graph on its \(n+1\) vertices is connected, so every vertex of
\(\sigma\) is within another \(n\) steps of that vertex.  Thus all
such simplices use vertices in the finite graph neighborhood
\[
 B^Y_{\lfloor s\sqrt n\rfloor+n}(A)
 :=\{v\in Y:d_Y(v,A)\leq\lfloor s\sqrt n\rfloor+n\}.
\]
There are only finitely many possibilities for \(\sigma\) and its
faces.  Likewise, every top simplex containing a specified vertex
\(v\) uses vertices in \(B^Y_n(\{v\})\); hence each vertex has a
finite star.

For completeness, the finite-subcomplex topology statement implies
the corresponding statement on all of \(|K|\).  Given \(p\), let
\(\eta=\min\{p(v):v\in\supp p\}>0\).  By
\eqref{eq:coordinate-control}, the ball of radius
\(\eta/(2\sqrt n)\) about \(p\) lies in the finite star of
\(\supp p\).  Within this star the metric and realization topologies
agree.  The open star of \(\supp p\) is also a neighborhood of \(p\)
in the locally finite realization, so the agreement is local in
both topologies.  A closed metric ball is a closed subset of the
finite union of closed simplices it meets, and is therefore compact.

If \(L_s(hc)\) meets \(\kappa\), \eqref{eq:support-distance} supplies
\(a\in A\) and \(b\in V(\kappa)\) with
\(d_Y(ha,b)\leq s\sqrt n\).  For each \(a\), the map \(h\mapsto ha\)
is injective because \(H\) acts freely on vertices.  The finitely
many pairs \(a,b\) and finite graph balls therefore allow only
finitely many \(h\).  Finally, for \(v\in V(\kappa)\), the same
injection maps \(H_\kappa\) into \(V(\kappa)\).  In particular,
\(|H_\kappa|\leq|V(\kappa)|\).
\end{proof}

We will also use the following consequence.  For fixed \(c\) and \(s\),
only finitely many \(h\in H\) satisfy
\begin{equation}
 L_s(hc)\cap L_s(c)\ne\varnothing.
 \label{eq:finite-ball-overlaps}
\end{equation}
Indeed, an intersection gives \(\rho(hc,c)\leq2s\), and then
\eqref{eq:support-distance} puts some vertex of \(h\supp c\)
within distance \(2s\sqrt n\) of \(\supp c\).  The same
injectivity argument applies.

\subsection{Simplex areas and distance gradients}

The metric now provides finite neighborhoods in which we can make
cuts.  For area comparisons, we use Euclidean structures separately
inside the top simplices.  They need not agree on shared faces.

Fix a top simplex \(\sigma\).  For each spanning tree \(T\) of the
connected graph on \(V(\sigma)\), equip its affine span with the
Euclidean structure in which the tree-edge vectors
\(\delta_v-\delta_w\), with either orientation for each edge, form
an orthonormal basis of the translation space.  For \(0\leq j\leq n\),
let \(a_j^\sigma\) be the average, over these finitely many trees,
of Euclidean \(j\)-area on polyhedral subsets of \(\sigma^\circ\)
of dimension at most \(j\).  We allow relatively open faces in
these sets; pieces of dimension less than \(j\) have zero
\(j\)-area, and \(a_0^\sigma\) is counting measure.  These area
functionals are covariant under \(H\), since the action permutes
the spanning trees.

\begin{lemma}\label{lem:tree}
In every tree Euclidean structure on a top simplex \(\sigma\), its
\(n\)-volume is \(1/n!\).  Every affine branch of \(f_c|_\sigma\)
has gradient of norm at most one, for every \(c\in|K|\).
The same gradient bound holds after restriction to any affine
subspace of the span of \(\sigma\).
\end{lemma}

\begin{proof}
Root the tree at \(v_0\), orient its edges away from the root, and
order the remaining vertices so that each parent precedes its
children.  Put
\[
 t_i=\delta_{v_i}-\delta_{\operatorname{parent}(v_i)},
 \qquad b_i=\delta_{v_i}-\delta_{v_0}
 \quad(1\leq i\leq n).
\]
Each \(b_i\) is the sum of the \(t_j\) along its root path.
Thus the change-of-basis matrix is triangular with diagonal
entries one.  The \(t_i\) are an orthonormal basis in the chosen
metric, so the parallelepiped spanned by the \(b_i\) has volume
one.  The simplex has volume \(1/n!\).

An affine branch has derivative along \(t_i\) equal, up to sign, to
\[
 \frac{\ell(v_i)-\ell(\operatorname{parent}(v_i))}{\sqrt n}
\]
for a \(1\)-Lipschitz potential \(\ell\).  Since the two vertices
are adjacent, this derivative has absolute value at most
\(1/\sqrt n\).  The squared gradient norm is the sum of the
squares of these \(n\) derivatives and is at most one.
Restriction to an affine subspace replaces the gradient by its
orthogonal projection, which cannot increase its norm.
\end{proof}

\subsection{The volume margin in high dimension}

We next bound the total top-dimensional area of a metric ball by
counting its possible connected vertex supports.  Fix
\[
 R=4D^2,\qquad r=R/2.
\]
For a specified vertex \(v\), there are at most \(D^{2n}\)
connected sets of \(n+1\) vertices containing \(v\).  To prove
this, choose a spanning tree of each such set rooted at \(v\).
A traversal using each tree edge once in each direction is a
walk of length \(2n\), starting at \(v\), whose visited set is
exactly the prescribed set.  Choosing one such walk per set
gives an injection into the length-\(2n\) walks from \(v\),
of which there are at most \(D^{2n}\).

Set \(k=\lfloor R\sqrt n\rfloor\).  A graph ball of radius \(k\)
has at most
\[
 \sum_{j=0}^kD^j\leq D^{k+1}
\]
vertices, since \(D\geq2\).  As \(|\supp c|\leq n+1\), there are
at most \((n+1)D^{k+1}\) vertices within distance \(k\) of
\(\supp c\).  By \eqref{eq:support-distance}, each top simplex
meeting \(L_R(c)\) contains one of them.  The connected-set
count and Lemma~\ref{lem:tree} therefore give, for every \(c\in|K|\),
\begin{equation}
 \sum_{\sigma\ \mathrm{top}}
 a_n^\sigma\bigl(L_R(c)\cap\sigma^\circ\bigr)
 \leq V_R:=
 \frac{(n+1)D^{\lfloor R\sqrt n\rfloor+1}D^{2n}}{n!}.
 \label{eq:volume-bound}
\end{equation}
Counting a simplex more than once only enlarges this upper bound.

With \(D\) fixed, the ratio of this estimate to
\((R-r)^n/n!\) is
\[
 \frac{V_R}{(R-r)^n/n!}
 =\frac{(n+1)D^{\lfloor4D^2\sqrt n\rfloor+1}}{2^n}
 \longrightarrow0.
\]
Indeed, its logarithm is
\(-n\log2+4D^2(\log D)\sqrt n+O_D(\log n)\).
We may therefore choose \(n_0(D)\geq1\) so that, for every
\(n\geq n_0(D)\),
\begin{equation}
 V_R<\frac{(R-r)^n}{n!}.
 \label{eq:strict-volume}
\end{equation}
For the remainder of the geometric proof, fix such an \(n\).
The graph degree bound \(D\), and hence \(R\) and \(r\), remain
fixed as dimensions increase.  The strict margin
\eqref{eq:strict-volume} will convert successive area comparisons
into a bound on the number of points at the bottom of a
separating filtration.
